% Matriz de consistencia simplificada (entregable semana 7): problema, objetivos, hipótesis,
% variables por hipótesis e indicadores. Versión completa: matriz_consistencia.tex.
% Compilar: latexmk -xelatex matriz_consistencia_simplificada.tex
\documentclass[10pt]{article}
\usepackage[a4paper,landscape,margin=1.5cm]{geometry}
\usepackage{fontspec}
\setmainfont{Arial}
\usepackage[spanish,es-noshorthands]{babel}
\usepackage{array}

\newcolumntype{L}[1]{>{\raggedright\arraybackslash}p{#1}}
\newcommand{\ingles}[1]{\textit{#1}}
\setlength{\parindent}{0pt}
\renewcommand{\arraystretch}{1.25}
\pagestyle{empty}

\begin{document}

{\centering
{\large\textbf{MATRIZ DE CONSISTENCIA}}\par\smallskip
Sistema de análisis satelital de inundaciones basado en modelos de segmentación profunda para el
mapeo de áreas inundadas en el Perú\par}
\bigskip

{\footnotesize\setlength{\tabcolsep}{4pt}
\begin{tabular}{|L{3.6cm}|L{4.6cm}|L{5.2cm}|L{5.6cm}|L{5.6cm}|}
\hline
\textbf{Problema} & \textbf{Objetivos} & \textbf{Hipótesis} &
\textbf{Variables (independientes y dependientes)} & \textbf{Indicadores} \\ \hline
¿De qué manera se puede desarrollar un sistema de análisis satelital de inundaciones para el
Perú, basado en modelos de segmentación profunda, que genere mapas de la extensión del agua y de
su incertidumbre a partir de imágenes Sentinel-2, aun cuando no se dispone de datos etiquetados
del Perú para entrenar los modelos?
&
\textbf{Objetivo general.} Desarrollar un sistema de análisis satelital de inundaciones para el
Perú, basado en modelos de segmentación profunda, que genere mapas de la extensión del agua y de
su incertidumbre a partir de imágenes Sentinel-2.
&
\textbf{Hipótesis general.} Los modelos fundacionales geoespaciales pierden menos desempeño que la
U-Net al pasar de Sen1Floods11 al Perú; la adaptación de dominio y la calibración mejoran en parte
la confiabilidad de los mapas, y el sistema cumple sus requisitos en un tiempo adecuado para la
respuesta después de un evento.
&
\textbf{Var. indep.:} modelo de segmentación, estrategia de entrenamiento o adaptación y
configuración del sistema.\newline
\textbf{Var. dep.:} calidad del mapeo de inundaciones.
&
IoU de agua; caída de IoU de Sen1Floods11 al Perú; ECE; AUROC; requisitos cumplidos (\%);
tiempo por escena (min).
\\ \cline{2-5}
&
\textbf{OE1.} Comparar el IoU de agua de TerraMind v1-B, Prithvi-EO-2.0 300M y U-Net en
Sen1Floods11 y en los eventos peruanos de 2023 y 2017, con tres esquemas de entrenamiento, para
seleccionar el modelo del sistema.
&
\textbf{HE1.} En Sen1Floods11 no se detecta una diferencia significativa de IoU de agua entre los
modelos fundacionales y la U-Net. En el Perú, la caída de IoU de los modelos fundacionales es
menor que la de la U-Net.
&
\textbf{Var. indep.:} arquitectura (TerraMind, Prithvi, U-Net) y esquema de entrenamiento
(Sen1Floods11; Perú; Sen1Floods11 y ajuste fino en el Perú).\newline
\textbf{Var. dep.:} exactitud de la segmentación y generalización al Perú.
&
IoU de agua; diferencia pareada de IoU con IC 95~\%; caída de IoU de Sen1Floods11 al Perú.
\\ \cline{2-5}
&
\textbf{OE2.} Determinar en qué medida la adaptación de dominio no supervisada y la calibración de
la incertidumbre mejoran la confiabilidad de los mapas en el Perú, e identificar las condiciones
en las que se concentran los errores.
&
\textbf{HE2.} Los errores se concentran en pocas condiciones (bordes de nubes, vegetación
inundada, terreno desértico o costero). La adaptación aumenta el IoU de agua sin llegar al modelo
entrenado en el Perú. El escalado de temperatura y los ensambles reducen el ECE, y la
incertidumbre permite detectar errores.
&
\textbf{Var. indep.:} método de adaptación (ninguno, igualación de histogramas, FDA, AdaBN) y
tratamiento de la incertidumbre (sin calibrar, escalado de temperatura, ensamble).\newline
\textbf{Var. dep.:} generalización al Perú, calibración y detección de errores.
&
Diferencia de IoU entre el modelo adaptado y el sin adaptar (IC 95~\%); proporción de errores
por condición; ECE (\%); AUROC de detección de errores ($>$ 0.5).
\\ \cline{2-5}
&
\textbf{OE3.} Implementar el sistema como una aplicación web y evaluar si cumple los requisitos
funcionales y el tiempo de procesamiento por escena Sentinel-2.
&
\textbf{HE3.} El sistema cumple el 100~\% de los requisitos funcionales RF1 a RF7 y procesa una
escena Sentinel-2 completa en 15 minutos o menos, incluida la descarga.
&
\textbf{Var. indep.:} configuración del sistema (modelo seleccionado, flujo de obtención y
procesamiento de imágenes).\newline
\textbf{Var. dep.:} desempeño del sistema.
&
Requisitos funcionales cumplidos (\%); tiempo de procesamiento por escena (min).
\\ \hline
\end{tabular}}

\medskip
{\footnotesize\raggedright
\textit{Nota.} IoU: intersección sobre unión; ECE: error de calibración esperado; AUROC: área bajo
la curva ROC; IC: intervalo de confianza. La versión completa, con diseño metodológico y
trazabilidad, está en \texttt{matriz\_consistencia.pdf}.\par\smallskip
Autores: Gerson Andre Garrido Mansilla, Renzo Elvis Cerron Tome. Elaboración: propia.}

\end{document}
